\section{Installing the residual vanishing theorem on the host frame}
\label{sec:hxp60}

Fix an arbitrary instance of the host problem. Thus $\Omega$ is the
specified algebraic closure of $\mathbb Q$, $\ell$ is a prime, $c$ is
the given coefficient datum, and $O=c.O$ has residue field $k$.
The supplied level $N\ne0$, exponent $e$, character
$\psi:(\mathbb Z/N)^\times\to O^\times$ and integer $o_1$ satisfy
the original $\mathrm{HostData}(\ell,e,\psi,o_1)$. Fix any $q$ satisfying
the host's $\mathrm{XData}(\ell,\psi,q)$. Write $G=\mathrm{Gal}(\Omega/\mathbb Q)$.
Take the supplied finite set $S$, its choices of valuation rings
$A_v$, the frame $F$, and all the identifications in
$\mathrm{IsHostFrame}$. In particular $S$ consists of the primes dividing
$\ell N$. The quotient $\Gamma=G_S$ kills the closed normal subgroup
generated by inertia above places outside $S$, and $F.\pi=\pi:G\to\Gamma$
is its quotient map. The distinguished $v_\ell$ lies over $\ell$.
Put $D_v=D_{A_v}$ and $I_v=I_{A_v}$. These are the actual subgroups
$F.D(v),F.I(v)$ of $G$, and $\mathrm{loc}_v$ is $\pi$ restricted to $D_v$.
The frame characters on $\Gamma$ are $\eta=F.\eta$ and $\Psi=F.\Psi$.
A bar on $\Psi$ denotes reduction $O^\times\to k^\times$.

% BEGIN HXP1
\paragraph{HXP1. The character and quotient for this frame.}
\label{claim:HXP1}
The scalar character needed for the weight $-2$ piece is
\[
   \lambda(s)=\bar\Psi(s)^{-2}\eta(s)\quad (s\in\Gamma).
\]
On $G$, $\gamma=\lambda\circ\pi$ is an inverse residual product
character covered by BR7 over $k$ with its given residue map.
The subgroup
\[
    K=\ker\big((\eta,\bar\Psi)\circ\pi\big)\trianglelefteq G
\]
lies in $\ker\gamma$ and has finite index prime to $\ell$.
Here and below an integer sign on the right of a field identity is
read in $k$ by its canonical map from $\mathbb Z$.

\begin{proof}
The ring $O$ is finite over $\mathbb Z_\ell$ and local with residue
characteristic $\ell$. Indeed the extension is integral, so
the contraction of its maximal ideal is the maximal ideal
$(\ell)$ of $\mathbb Z_\ell$. Consequently $k$ is a field quotient
of the finite-dimensional $\mathbb F_\ell$-algebra $O/\ell O$.
Write $|k|=\ell^r$ with $r\ge1$. The group $\eta$ takes values in
$\mathbb Z^\times=\{1,-1\}$, before reduction.

HXD3, applied with precisely the fixed $c,N,e,\psi,o_1,q$ and
$\Omega$, gives $d,f,n$, $\varphi$ of level $f$, and $\xi$ of level
$2^n$ over $O$. Their full Bernoulli character datum holds and
$f2^n\ne0$. For every $g\in G$ the residue of the Galois character of
$\operatorname{prodX}(\varphi,\xi)$ equals
\[
    \overline{\operatorname{Chr.Psi}(\ell,e,\operatorname{dirG}\psi)(g)}^{\,2}
       \,\eta_4(g)
      =\bar\Psi(\pi g)^2\eta(\pi g).
\]
The last equality uses both character identities of this host frame.
Since an integer sign equals its inverse, inverting this identity gives
$\gamma$ pointwise. We can therefore select the inverse alternative in
BR7 with $\mathrm{red}:O\to k$ the original residue map.
The signs and the residual units together have at most
$2(\ell^r-1)$ elements; more precisely the order of the image of
$(\eta,\bar\Psi)\circ\pi$ divides this number by Lagrange's theorem.
The host has $\ell\ne2$, hence the index of its kernel $K$ is prime
to $\ell$. Both character factors are one on $K$, which gives
$K\le\ker\gamma$.
\end{proof}
% END HXP1

% BEGIN HXP2
\paragraph{HXP2. The away-from-$\ell$ condition at every valuation.}
\label{claim:HXP2}
Let $z:\Gamma\to k$ be a locally constant $\lambda$-cocycle, and put
$u=z\circ\pi$. The frame's specified subgroup at $v$ is
\[
 N_v=\{s\in D_v : s\in I_v,\ \eta(\pi s)=1,\
                                     \bar\Psi(\pi s)=1\}.
\]
Suppose $z(\mathrm{loc}_v s)=0$ on $N_v$ at every $v\ne v_\ell$.
Then for every valuation ring $A$ of $\Omega$ with
$\operatorname{val}_A(\ell)\not<1$, the restriction of $u$ to
$I_A\cap K$ vanishes.

\begin{proof}
The quotient map is continuous, so $u$ is locally constant and
$u(gh)=u(g)+\gamma(g)u(h)$. In particular $u(1)=0$. For
$x\in\ker\gamma$, calculate
\[
   u(\sigma^{-1}x\sigma)
     =u(\sigma^{-1})+\gamma(\sigma)^{-1}
            \big(u(x)+u(\sigma)\big)
     =\gamma(\sigma)^{-1}u(x).                         \tag{HXP.1}
\]
This identity permits comparison with the one valuation above
each rational place chosen by the frame.

The valuation ring $A=\Omega$ has identity residue-field map, hence its
inertia is trivial. Every other $A$ lies above a rational prime $p$.
Indeed $A$ contains all algebraic integers, and its maximal ideal
contracts to a nonzero prime of $\mathbb Z$. If the contraction were
zero, every nonzero integer would be a unit in $A$; clearing an
integer denominator of an arbitrary algebraic number would then put
that number in $A$, forcing $A=\Omega$. The prime of contraction
determines the unique $v$ of $\mathbb Q$ below $A$. The hypothesis on
$\ell$ says $p\ne\ell$.

If $v\notin S$, the image of $I_A$ in $\Gamma$ is trivial by the
definition of $G_S$. Hence $u$ vanishes even on all of $I_A$.
If $v\in S$, the place differs from $v_\ell$. Conjugacy of valuations
above $v$ gives $\sigma\in G$ carrying $A_v$ to $A$. To recall
the infinite-extension point, on each finite normal subextension the
Galois group is transitive on the primes above $p$. The cosets of
lifts taking the first prime to the second are nonempty closed
subsets of the profinite $G$ and have the finite-intersection
property by passing to the compositum. An element in the intersection
carries the valuation rings to one another and conjugates their
inertia groups. For $x\in I_A\cap K$, it follows that
$s=\sigma^{-1}x\sigma\in I_{A_v}\cap K$ by normality of $K$.
Regarded as an element of $D_v$, $s$ belongs to the displayed $N_v$,
so $u(s)=z(\mathrm{loc}_v s)=0$. Equation (HXP.1) gives $u(x)=0$
since $K\le\ker\gamma$ and $\gamma(\sigma)$ is a unit. This also
handles every place in the host's special set away from $\ell$.
\end{proof}
% END HXP2

% BEGIN HXP3
\paragraph{HXP3. The exact tame-index condition and global scalar.}
\label{claim:HXP3}
Retain HXP2's local conditions on $z$ and suppose at the
distinguished place there exists $m_\ell\in k$ such that
\[
 z(\mathrm{loc}_{v_\ell}s)=\lambda(\mathrm{loc}_{v_\ell}s)m_\ell-m_\ell
                  \qquad(s\in D_{v_\ell}).
\]
Then some $m\in k$ satisfies
$z(s)=\lambda(s)m-m$ for every $s\in\Gamma$.

\begin{proof}
To apply BR7 on $G$, its original condition (BR.2) asks at every
valuation ring $A$ for a subgroup $H_A\le G$. With
$J_A=I_A\cap\ker\gamma$, the index $[J_A:J_A\cap H_A]$ must be
finite and prime to $\ell$, and $u$ must vanish on $J_A\cap H_A$.
If $\operatorname{val}_A(\ell)\not<1$, take $H_A=K$. The quotient
$J_A/(J_A\cap K)$ injects into $G/K$. Its order divides the
prime-to-$\ell$ order from HXP1. HXP2 supplies the required
vanishing, including for the trivial valuation.

Otherwise $A$ lies over $\ell$, as does $A_{v_\ell}$ by the
host's place equality. Choose $\sigma$ carrying $A_{v_\ell}$ to
$A$ as in HXP2. For $x\in J_A$, the element
$s=\sigma^{-1}x\sigma$ belongs to $I_{A_{v_\ell}}\subset D_{v_\ell}$
and $\gamma(s)=1$. The assumed local coboundary therefore gives
$u(s)=0$. Equation (HXP.1) gives $u(x)=0$. We may take $H_A=G$,
with relative index one. This proves (BR.2) for every $A$.
BR7 for the inverse character in HXP1 gives $m\in k$ with
$u(g)=\gamma(g)m-m$ for every $g\in G$. The quotient map is
surjective, and both sides are pullbacks from $\Gamma$. Evaluating at
any lift of $s\in\Gamma$ gives the asserted global formula there.
\end{proof}
% END HXP3

% BEGIN HXP4
\paragraph{HXP4. The weight $-2$ Selmer quotient on the CountV host.}
\label{claim:HXP4}
For every supplied datum and frame fixed above,
\[
    F.\mathrm{selmerCohomologyQuotient}(-2)=0 .
\]
Here zero means that every class is the zero class. This is
$F.\mathrm{Xpp}$, and the universal statement is precisely
$\mathrm{HostXpp}$.

\begin{proof}
Index the three rows and columns by $0,1,2$. The diagonal weights
are $1,0,-1$, so the matrix space $M$ of pure weight $-2$ consists
of the matrices $aE_{20}$, $a\in k$. In the base representation,
diagonal conjugation multiplies entry $(i,j)$ by
$\bar\Psi(s)^{\mathrm{wt}(i)-\mathrm{wt}(j)}$. When
$\eta(s)=-1$ the adjoint action first sends $X$ to $-JX^tJ$,
where $J$ reverses the three indices. The $(2,0)$ entry of this
new matrix is $-X_{20}$, because reversal and transposition fix
that position together. When $\eta(s)=1$ this step is the identity.
It follows that on $M$ the residual adjoint action is exactly
\[
    F.\bar g(-2)(s)(aE_{20})
              =\bar\Psi(s)^{-2}\eta(s)\,aE_{20}=\lambda(s)aE_{20}.
                                                               \tag{HXP.2}
\]

Let $b:\Gamma\to M$ represent any element of the actual Selmer
quotient. Its numerator consists of continuous cocycles whose
restriction to each $D_v$ belongs to the prescribed local submodule
$F.ZL(-2,v)$. Its denominator consists of those members of this
numerator which are global coboundaries. Since $M$ is discrete,
$z(s)=b(s)_{20}$ is locally constant. Taking that entry in the
cocycle identity, (HXP.2) gives
$z(st)=z(s)+\lambda(s)z(t)$. For $v\ne v_\ell$ the local submodule
requires $b(\mathrm{loc}_v g)=0$ for $g\in N_v$, so the scalar
equation of HXP2 holds. At $v_\ell$ the weight is strictly negative,
so the restriction of $b$ to the entire $D_{v_\ell}$ is a
coboundary. Taking the $(2,0)$ entry of the local witness gives
precisely the $m_\ell$ required in HXP3. The additional clause
in the local definition is indexed by weight zero and has false
antecedent here.

Apply HXP3 and set $B=mE_{20}\in M$. For every $s\in\Gamma$,
the matrix $F.\bar g(-2)(s)B-B$ and $b(s)$ have the same $(2,0)$
entry by (HXP.2); all their other entries are zero.
Hence $b$ belongs to the indicated denominator. This proves the
claimed vanishing for the original quotient over $k$.

In the CountV application a host $h$ supplies $o_1=h.o$, $q=h.q$
and $N=h.N$. The source extracts the two arithmetic predicates from
the character of this host, then chooses $S,A,F$ and the witness
of $\mathrm{IsHostFrame}$ once. Substituting exactly those objects
into the proved universal statement supplies the $F.\mathrm{Xpp}$
argument of $\mathrm{abstract\_count}$. The chosen marked element,
universal representation and length estimates are other arguments of
that application. This identifies the input at the original call site
without adding any hypotheses to $\mathrm{HostXpp}$.
\end{proof}
% END HXP4
